\documentclass[UTF8,11pt,a4paper]{ctexart}
\usepackage[a4paper,left=1.45cm,right=1.45cm,top=1.55cm,bottom=1.65cm]{geometry}
\usepackage{amsmath,amssymb,bm}
\usepackage{graphicx}
\usepackage{booktabs,tabularx,array,longtable}
\usepackage{enumitem}
\usepackage[dvipsnames,svgnames]{xcolor}
\usepackage{tikz}
\usetikzlibrary{arrows.meta,positioning,calc,fit,shapes.geometric}
\usepackage[most]{tcolorbox}
\usepackage{hyperref}
\usepackage{caption}
\usepackage{float}
\usepackage{fancyhdr}
\usepackage{titlesec}
\usepackage{microtype}
\IfFontExistsTF{Times New Roman}{\setmainfont{Times New Roman}}{\setmainfont{Liberation Serif}}
\IfFontExistsTF{Noto Sans CJK SC}{\setCJKmainfont{Noto Sans CJK SC}}{}
\IfFontExistsTF{Noto Sans CJK SC}{\setCJKsansfont{Noto Sans CJK SC}}{}
\IfFontExistsTF{DejaVu Sans Mono}{\setmonofont{DejaVu Sans Mono}}{\setmonofont{Latin Modern Mono}}
\definecolor{mainblue}{HTML}{0B5CAD}
\definecolor{softblue}{HTML}{EAF4FF}
\definecolor{darktext}{HTML}{1F2937}
\definecolor{softgray}{HTML}{F5F7FA}
\definecolor{accent}{HTML}{F59E0B}
\definecolor{goodgreen}{HTML}{2E7D32}
\definecolor{badred}{HTML}{B91C1C}
\definecolor{purple}{HTML}{6D5BD0}
\hypersetup{colorlinks=true,linkcolor=mainblue,urlcolor=DarkOrange,citecolor=mainblue}
\setlength{\parindent}{2em}
\setlength{\headheight}{14pt}
\setlength{\parskip}{0.22em}
\setlist[itemize]{leftmargin=2em,itemsep=0.23em,topsep=0.25em}
\setlist[enumerate]{leftmargin=2.2em,itemsep=0.23em,topsep=0.25em}
\renewcommand{\arraystretch}{1.2}
\captionsetup{font=small,labelfont=bf}
\titleformat{\section}{\Large\bfseries\sffamily\color{mainblue}}{\thesection}{0.7em}{}
\titleformat{\subsection}{\large\bfseries\sffamily\color{darktext}}{\thesubsection}{0.6em}{}
\titleformat{\subsubsection}{\normalsize\bfseries\sffamily\color{darktext}}{\thesubsubsection}{0.55em}{}
\pagestyle{fancy}
\fancyhf{}
\fancyhead[L]{\small\sffamily Exploring Diffusion Transformer Designs via G}
\fancyhead[R]{\small\sffamily \leftmark}
\fancyfoot[C]{\small\thepage}
\renewcommand{\headrulewidth}{0.4pt}
\newtcolorbox{infobox}[1]{enhanced,breakable,colback=softblue,colframe=mainblue!70!black,boxrule=0.8pt,arc=2mm,left=1.5mm,right=1.5mm,top=1mm,bottom=1mm,title=\bfseries #1}
\newtcolorbox{keybox}[1]{enhanced,breakable,colback=softgray,colframe=darktext!45,boxrule=0.6pt,arc=2mm,left=1.5mm,right=1.5mm,top=1mm,bottom=1mm,title=\bfseries #1}
\newtcolorbox{notebox}[1]{enhanced,breakable,colback=accent!8,colframe=accent!80!black,boxrule=0.7pt,arc=2mm,left=1.5mm,right=1.5mm,top=1mm,bottom=1mm,title=\bfseries #1}
\newcommand{\term}[1]{\textbf{\textcolor{mainblue}{#1}}}
\begin{document}

\section{论文基本信息}
\begin{center}
\begin{tabularx}{0.94\linewidth}{>{\bfseries}p{0.18\linewidth}X}
\toprule
题目 & Exploring Diffusion Transformer Designs via Grafting\\
课件日期 & 2025-06-19\\
研究方向 & 扩散 Transformer / 架构编辑\\
一句话概括 & 把新模块“嫁接”到已有预训练 DiT 中,用低成本局部训练评估架构设计,避免每次创新都从头完整训练。\\
\bottomrule
\end{tabularx}
\end{center}

\begin{infobox}{摘要要点}
把新模块“嫁接”到已有预训练 DiT 中,用低成本局部训练评估架构设计,避免每次创新都从头完整训练。 本说明按“研究问题—核心机制—基础公式—实验阅读—局限与优化”的顺序重新组织，不保留原始材料的封面、目录和页码对应。
\end{infobox}

\section{研究问题与动机}
这项工作的出发点不是简单地替换某个网络层，而是识别现有范式中最影响\term{质量、效率、可靠性或可部署性}的瓶颈。结合论文所讨论的问题，可以将研究动机概括为以下几点：
\begin{itemize}
\item 应用价值 局限性与未来工作 Exploring Diffusion Transformer Designs via Grafting
\item ·计算特性: 并行块提升SM利用率(论文6章) 减少层间同步点 ·实测数据: 2倍理论速度提升(Table 6备注) 适合大batch推理 ·内存优化: 减少中间激活存储 更均衡的显存占用 Exploring Diffusion Transformer Designs via Grafting
\item 局限性与未来工作 Exploring Diffusion Transformer Designs via Grafting ·现有局限: 仅处理MHA/MLP(论文8章) 需预训练模型 ·技术瓶颈: 合成数据质量影响(Figure D.4) 深度压缩的极限 ·未来方向: 自动化嫁接策略 跨模态架构移植
\item Exploring Diffusion Transformer Designs via Grafting 架构创新=从头训练?这个认知该打破了!
\item 传统流程:新想法→完整训练→评估(耗时耗力) 类比说明:相当于每次改装汽车都要重造发动机 论文突破:对预训练模型进行编辑(论文第3章)
\item 为什么GPU更爱"宽扁"架构?
\end{itemize}

\begin{notebox}{理解重点}
阅读这类论文时，应把“问题是否真实存在”“方法是否直接解决该问题”“实验是否排除了其他解释”分开判断。仅凭更大的模型、更多数据或更长训练得到的提升，不能自动视为结构创新带来的收益。
\end{notebox}

\section{核心思想与整体流程}
把新模块“嫁接”到已有预训练 DiT 中,用低成本局部训练评估架构设计,避免每次创新都从头完整训练。 从信息流角度看，其基本流程可以整理为下图。

\begin{figure}[H]
\centering
\begin{tikzpicture}[>=Latex,node distance=0.78cm]
\tikzstyle{procstep}=[draw,rounded corners,fill=softblue,align=center,minimum width=3.05cm,minimum height=0.95cm,text width=2.75cm]
\node[procstep] (a) {预训练 DiT};
\node[procstep,right=of a] (b) {模块替换/嫁接};
\node[procstep,right=of b] (c) {短程恢复训练};
\node[procstep,right=of c] (d) {架构评估};
\draw[->,thick,mainblue] (a)--(b);
\draw[->,thick,mainblue] (b)--(c);
\draw[->,thick,mainblue] (c)--(d);
\end{tikzpicture}
\caption{核心信息流与处理阶段。}
\end{figure}

\subsection{方法组成与信息流}
\begin{itemize}
\item Exploring Diffusion Transformer Designs via Grafting 架构创新=从头训练?这个认知该打破了!
\item 传统流程:新想法→完整训练→评估(耗时耗力) 类比说明:相当于每次改装汽车都要重造发动机 论文突破:对预训练模型进行编辑(论文第3章)
\item 应用价值 局限性与未来工作 Exploring Diffusion Transformer Designs via Grafting
\item 生物嫁接 × AI模型 = 突破性创新 类比解析: 砧木:预训练模型(Baseline) 接穗:新架构算子(Hyena-X/Y等,论文4.1节) 形成层:激活蒸馏 关键差异:植物嫁接(南瓜砧木)成功率≈95\%,模型嫁接FID变化仅+0.37(Table 4最佳案例) Exploring Diffusion Transformer Designs via Grafting
\item 局限性与未来工作 Exploring Diffusion Transformer Designs via Grafting ·现有局限: 仅处理MHA/MLP(论文8章) 需预训练模型 ·技术瓶颈: 合成数据质量影响(Figure D.4) 深度压缩的极限 ·未来方向: 自动化嫁接策略 跨模态架构移植
\item Exploring Diffusion Transformer Designs via Grafting 从Stable Diffusion到Sora:扩散模型的爆发与瓶颈 行业现状:扩散模型已成为AIGC核心(图像/视频生成) 核心矛盾:模型越强→训练成本越高(DiT-XL/2需7000K迭代) 从头开始训练模型的成本过高 数据支撑:Table 2显示完整训练需100\%,嫁接仅需0.63\%-10\% training data
\item 9\%(Table 4) 替换策略: 交错替换优于全替换(50\% vs 100\%对比) 局部性强的层优先(Figure 2指导) Exploring Diffusion Transformer Designs via Grafting Attention Is All You Need
\end{itemize}

\begin{keybox}{结构解析}
\begin{itemize}
\item 2 从Stable Diffusion到Sora：扩散模型的爆发与瓶颈 行业现状：扩散模型已成为AIGC核心（图像/视频生成） 核心矛盾：模型越强→训练成本越高（DiT-XL/2需7000K迭代） 从头开始训练模型的成本过高 数据支撑：Table 2显示完整训练需……
\end{itemize}
\end{keybox}

下图补充展示模块连接、数据流或训练阶段之间的关系。
\begin{figure}[H]
\centering
\includegraphics[width=0.94\linewidth,height=0.43\textheight,keepaspectratio]{method_diagram.jpg}
\caption{模型结构、数据流或主要训练阶段。}
\end{figure}

\section{数学与机制理解}
本节给出理解该工作的基础数学结构。公式用于解释方法所属范式，并不替代原论文中的完整推导。
\begin{equation}
x_t=\alpha_t x_0+\sigma_t\varepsilon,\qquad \varepsilon\sim\mathcal N(0,I)
\end{equation}
扩散类方法把真实样本与噪声按时间混合，模型再学习逆向去噪、速度场或条件流。不同论文的主要差异通常落在表示空间、网络骨干和训练目标上。

\begin{keybox}{从公式回到方法}
应重点观察论文究竟改变了公式中的哪一部分：是输入表示、连接矩阵、条件变量、参数化方式、训练目标，还是求解与调度过程。真正的创新通常可以在这一层被准确定位。
\end{keybox}


\section{实验结果应如何阅读}
\begin{itemize}
\item 9\%(Table 4) 替换策略: 交错替换优于全替换(50\% vs 100\%对比) 局部性强的层优先(Figure 2指导) Exploring Diffusion Transformer Designs via Grafting Attention Is All You Need
\item 文生图实战 方案: 替换层:7-27层MHA(论文5章) 新算子:Hyena-X(K=4) 训练数据:12k合成样本 结果: 速度↑:1.43倍(235.46ms→164.58ms) 质量↓:GenEval
\item 75(Table 5) 特殊技巧: 使用LoRA(rank=64) 分层clip norm(0.1/0.01) Exploring Diffusion Transformer Designs via Grafting
\item 77超越同类深度模型(对比Table 6中TinyDiT-D14) 参数仅增6\%(712M vs 675M) 硬件优势: 更适合GPU并行计算 减少层间通信开销 Exploring Diffusion Transformer Designs via Grafting
\item 77 vs 剪枝2.86(Table 6) 从头训练需要7000k次迭代(论文2章) 成本对比: 嫁接:10\%数据+230K步 剪枝:100\%数据+500K步(TinyDiT) 灵活性对比: 嫁接支持架构重组 剪枝仅能减少参数 Exploring Diffusion Transformer Designs via Grafting
\item ·计算特性: 并行块提升SM利用率(论文6章) 减少层间同步点 ·实测数据: 2倍理论速度提升(Table 6备注) 适合大batch推理 ·内存优化: 减少中间激活存储 更均衡的显存占用 Exploring Diffusion Transformer Designs via Grafting
\item 深度压缩 技术原理: 并行化相邻块(论文6章"parallelizing every pair of blocks") 线性合并输出(Figure 5(b)结构) 性能表现: FID
\end{itemize}

\begin{notebox}{结果解读}
\begin{itemize}
\item 7 注意力机制的替代方案 候选方案： Hyena-X/Y（论文原创，K=4） 滑动窗口注意力（w=4） Mamba-2（ds=64） 性能对比： 最佳FID：Hyena-Y 2.72（vs 基线2.27） 最大FLOPs↓：Hyena-X 49.9\\%（Table ……
\end{itemize}
\end{notebox}

该部分以文字概括实验结论与比较条件，避免用整页截图代替分析。
实验部分不应只看单一最好数字。还需要核对模型规模、训练数据、训练步数、采样或解码预算、硬件条件以及是否使用额外蒸馏、检索或后处理。只有比较条件接近时，结果才能真正说明方法本身的优势。

\section{优点、局限与可优化空间}
\subsection{主要优点}
\begin{itemize}
\item \term{问题与方法对应明确}：论文围绕一个可识别的核心瓶颈展开，方法结构能够直接映射到该瓶颈。
\item \term{可复用的设计思想}：即使不完整复现模型，其中的分解、稀疏化、递归、条件控制或系统优化思想也可迁移到相邻任务。
\item \term{实验提供了多角度证据}：实验通常同时展示主结果、效率比较、案例或消融，使结论不完全依赖单一指标。
\end{itemize}

\subsection{局限与进一步优化}
\begin{itemize}
\item 需要进一步区分\term{结构收益}与更大算力、更多数据、不同训练技巧带来的收益，并在等参数、等计算预算下进行严格比较。
\item 对超参数、输入分布和模型规模的敏感性仍应系统分析；在一个数据集上的最优配置不一定能直接迁移。
\item 实际部署还要考虑显存峰值、并发、延迟抖动、失败案例和鲁棒性，而不仅是离线平均指标。
\item 可尝试将该方法与蒸馏、量化、缓存复用、动态路由或更强的表示学习结合，验证不同优化是否互补。
\end{itemize}

\section{结论}
总体而言，把新模块“嫁接”到已有预训练 DiT 中，用低成本局部训练评估架构设计，避免每次创新都从头完整训练。。 进一步需要注意：局限性与未来工作 Exploring Diffusion Transformer Designs via Grafting ·现有局限: 仅处理MHA/MLP(论文8章) 需预训练模型 ·技术瓶颈: 合成数据质量影响(Figure D.4) 深度压缩的极限 ·未来方向: 自动化嫁接策略 跨模态架构移植;应用价值 局限性与未来工作 Exploring Diffusion Transformer Designs via Grafting。

\end{document}
